\section{Results}
\label{sec:results}

\subsection{Main Result: 5$\times$ Drift Slope Difference}

Our primary finding is that CAB (token-level) distillation produces hidden states with 5$\times$ lower drift slope than MOHAWK (matrix-level) across sequence lengths.

\textbf{Quantitative results} (H-M2):

\begin{table}[h]
\centering
\begin{tabular}{llll}
\toprule
\textbf{Metric} & \textbf{CAB} & \textbf{MOHAWK} & \textbf{Ratio} \\
\midrule
Drift slope & 0.00090506 & 0.00452495 & 5.0$\times$ \\
Drift ratio (2048/512) & 1.34 & 2.02 & --- \\
95\% CI (slope) & [0.0008, 0.0010] & [0.0040, 0.0050] & --- \\
\bottomrule
\end{tabular}
\caption{Comparison of drift metrics between CAB and MOHAWK distillation objectives.}
\label{tab:drift_results}
\end{table}

The slope difference is statistically significant (p < 0.001). Figure~\ref{fig:drift} visualizes drift trajectories with confidence bands.

\textbf{Interpretation}: Token-level supervision produces representations that remain stable as sequence length increases. Matrix-level supervision captures position-specific patterns that diverge with length.

\subsection{Teacher Context Limit Discovery}

H-M1 revealed an unexpected architectural constraint: Phi-1.5 cannot process sequences beyond 2048 tokens.

\textbf{Observation}: At position 2049, Phi-1.5 raises IndexError on the position embedding matrix. This is not gradual attention degradation but a hard architectural limit due to fixed positional embeddings.

\textbf{Impact}: This constrains our experimental scope to $\leq$2048 tokens. True length extrapolation experiments (16K--32K as originally hypothesized) require teachers with relative position encodings (RoPE, ALiBi).

\subsection{Bounded vs Unbounded Drift}

Beyond slope, we observe qualitatively different drift behavior:

\textbf{CAB}: Drift ratio 1.34 indicates bounded degradation. Hidden state similarity at 2048 tokens remains within 34\% of the 512-token baseline. Across all measured layers (8, 12, 16), CAB maintains this bounded pattern.

\textbf{MOHAWK}: Drift ratio is 2.02 (computed as 13.79/6.84 from raw drift values). Per-layer analysis (Figure~\ref{fig:heatmap}) shows middle layers (12, 16) exhibit greatest divergence, consistent with attention patterns becoming more length-dependent in deeper layers.

\subsection{Summary of Hypothesis Outcomes}

\begin{table}[h]
\centering
\begin{tabular}{llll}
\toprule
\textbf{Hypothesis} & \textbf{Gate} & \textbf{Result} & \textbf{Evidence} \\
\midrule
H-E1 & MUST\_WORK & \textbf{PASS} & Both objectives trainable \\
H-M1 & MUST\_WORK & \textbf{PASS} & 2048 limit confirmed \\
H-M2 & MUST\_WORK & \textbf{PASS} & 5$\times$ slope difference \\
H-M3 & MUST\_WORK & \textbf{FAIL} & p=0.881 (PoC limitation) \\
\bottomrule
\end{tabular}
\caption{Summary of hypothesis outcomes.}
\label{tab:outcomes}
\end{table}

Three of four hypotheses pass. H-M3 was \textbf{not supported} (p=0.881 indicates no detectable interaction effect). This null result may reflect insufficient training data in PoC mode or a genuine absence of interaction between objective type and F1 retention.
